\section{Tessuto epiteliale}

\subsection{I tessuti: livelli di organizzazione}

\textbf{Tessuto}: insieme di cellule specializzate e di prodotti cellulari che esegue un numero di funzioni relativamente limitato, ma altamente specifico.

\begin{itemize}
  \item \textbf{epiteliale}: riveste le superfici esterne ed interne del corpo, le vie interne di transito di fluidi, e costituisce le ghiandole;
  \item \textbf{connettivo}: fornisce sostegno strutturale, riempie gli spazi tra gli organi, trasporta materiali ed immagazzina energia;
  \item \textbf{muscolare}: permette la contrazione di muscoli scheletrici, viscerali e del cuore;
  \item \textbf{nervoso}: trasmette ed integra impulsi elettrici da una parte all'altra del corpo.
\end{itemize}

\rubrica{Livelli di organizzazione}
$molecole \rightarrow cellule \rightarrow tessuti \rightarrow organi \rightarrow apparati$
\begin{itemize}
  \item \textbf{livello chimico o molecolare}: atomi che si uniscono a formare molecole con forme complesse (la forma specializzata di una molecola ne determina la funzione); studiato dalla biologia molecolare (anatomia sub-microscopica);
  \item \textbf{livello cellulare}: molecole che si uniscono a formare organuli con specifiche funzioni, componenti strutturali e funzionali delle cellule, le più piccole unità viventi del corpo; studiato dall'anatomia microscopica-citologia;
  \item \textbf{livello tissutale}: gruppo di cellule specializzate, anche di tipo diverso tra loro, che cooperano strutturalmente per svolgere una o più funzioni specifiche; studiato dall'anatomia microscopica-istologia;
  \item \textbf{livello degli organi}: due o più tessuti strutturalmente connessi per svolgere svariate funzioni (es.\ stomaco, cuore, cervello); a questo livello si pone il confine tra anatomia macroscopica e anatomia microscopica;
  \item \textbf{livello d'apparato}: più organi coordinati per una funzione comune (es.\ apparato cardiocircolatorio).
\end{itemize}

\rubrica{Differenziazione cellulare}
Ciascuna cellula riduce e limita selettivamente alcune proprie capacità operative e ne amplifica altre, aumentando l'efficacia delle funzioni rimaste. Produce nel corpo umano circa 200 tipi diversi di cellule specializzate; la combinazione di diversi tipi cellulari specializzati dà origine ai tessuti.

% ---------------------------------------------------------------------------
\subsection{Tessuto epiteliale: caratteristiche generali}

Il tessuto epiteliale si distingue in:
\begin{itemize}
  \item \textbf{epitelio di rivestimento}: riveste cute, cavità interne, superfici interne di arterie e vene, canale alimentare, vie respiratorie;
  \item \textbf{epitelio ghiandolare}: strutture che producono secrezioni liquide e/o viscose.
\end{itemize}

\rubrica{Caratteristiche}
\begin{itemize}
  \item \textbf{cellularità}: 100\% (quasi totale assenza di sostanza extracellulare);
  \item \textbf{polarità}: 100\% (organizzazione strutturale e funzionale orientata);
  \item \textbf{posizione delle cellule}: presenza di membrana basale su cui sono poste tutte le cellule;
  \item \textbf{vascolarizzazione}: 0\% (le cellule vengono nutrite mediante diffusione dei nutrienti dal liquido extracellulare);
  \item \textbf{innervazione}: 100\%;
  \item \textbf{proliferazione}: 100\%, per assicurare il ricambio cellulare (l'epitelio è un tessuto \textit{labile}/\textit{stabile}).
\end{itemize}

Le cellule epiteliali hanno forma poliedrica ed entrano in rapporto da un lato con un lume o una cavità (interna o esterna) e dall'altro con un tessuto connettivo. La \textbf{lamina basale} è costituita da materiale amorfo extracellulare a cui le superfici basali delle cellule epiteliali aderiscono mediante \textbf{emidesmosomi}; consente alle cellule epiteliali di mantenersi meccanicamente adese al tessuto connettivo. È composta da:
\begin{itemize}
  \item \textbf{lamina lucida};
  \item \textbf{lamina densa};
\end{itemize}
insieme costituiscono la lamina basale propriamente detta; sotto di essa si trova la \textbf{lamina reticolare} (o lamina fibroreticolare), che la connette al connettivo sottostante.

% ---------------------------------------------------------------------------
\subsection{Funzioni}

\begin{itemize}
  \item \textbf{protezione fisica}: rivestimento contro abrasioni, essiccamento, distruzione da agenti lesivi (es.\ ulcera);
  \item \textbf{permeabilità ed assorbimento}: di acqua, ioni, sostanze selezionate;
  \item \textbf{secrezione}: produzione di sostanze (\textit{secreti}), scaricate all'esterno delle cellule epiteliali secernenti (organizzate in strutture dette ghiandole, singole o in gruppi), con molteplici funzioni utili all'organismo.
\end{itemize}

In sintesi: protezione e barriera con captazione di stimoli di varia natura; secrezione, escrezione e assorbimento; scambio e trasporto di sostanze.

% ---------------------------------------------------------------------------
\subsection{Classificazione morfologica}

La morfologia è strettamente correlata alla funzione. Si classifica in base al numero di strati di cellule e alla forma delle cellule dello strato superficiale:

\begin{center}
\begin{forest} classalbero
[Epitelio
  [Monostratificato (singolo strato)
    [Pavimentoso semplice (cellula piatta)]
    [Cubico semplice (cellula cubica)]
    [Cilindrico semplice (cellula prismatica)]
  ]
  [Pluristratificato (multistrato)
    [Pavimentoso stratificato (cellula piatta)]
    [Cubico stratificato (cellula cubica)]
    [Cilindrico stratificato (cellula prismatica)]
  ]
]
\end{forest}
\end{center}

Esistono inoltre due varianti particolari, non riconducibili direttamente allo schema semplice/stratificato per numero di strati:
\begin{itemize}
  \item \textbf{epitelio pseudostratificato}: monostrato in cui i nuclei, a diverse altezze, danno la falsa impressione di più strati;
  \item \textbf{epitelio di transizione}: multistrato di cellule cubiche, capace di variare forma (vedi \S\ref{ssec:transizione}).
\end{itemize}

Gli epiteli semplici o monostratificati consentono assorbimento, diffusione e secrezione; gli epiteli stratificati garantiscono protezione da sollecitazioni meccaniche, fisiche e chimiche.

% ---------------------------------------------------------------------------
\subsection{Epiteli semplici (monostratificati)}

\rubrica{Epitelio pavimentoso semplice}
Cellule piatte in unico strato.
\begin{itemize}
  \item \textbf{alveolo polmonare}: permette scambi di sostanze nutritizie, sostanze di rifiuto e gas;
  \item \textbf{endotelio} di cuore e vasi sanguigni;
  \item \textbf{mesotelio} di cavità pleuriche ed addominale;
  \item foglietto parietale della capsula di Bowman, che circonda il glomerulo renale.
\end{itemize}

\rubrica{Epitelio cubico semplice}
Cellule cubiche in unico strato, con nucleo tondeggiante posto al centro della cellula.
\begin{itemize}
  \item \textbf{intestino tenue};
  \item \textbf{ghiandole} e loro dotti escretori;
  \item \textbf{tubuli renali}: secerne e riassorbe acqua e piccole molecole.
\end{itemize}

\rubrica{Epitelio cilindrico semplice}
Cellule prismatiche in unico strato.
\begin{itemize}
  \item riveste la maggior parte del \textbf{tubo digerente} (ricco di \textit{cellule caliciformi mucipare}, il cui citoplasma appare scarsamente colorabile perché ricco di \textit{mucinogeno}, che protegge la mucosa dall'acidità gastrica);
  \item riveste \textbf{cistifellea} e dotti escretori di ghiandole esocrine.
\end{itemize}

% ---------------------------------------------------------------------------
\subsection{Epitelio pseudostratificato}

Monostrato in cui le cellule, di altezza diseguale, hanno i nuclei posti a diverse quote, dando la falsa impressione di uno stratificato. Esempio: epitelio pseudostratificato ciliato della trachea, tra le cui cellule sono presenti numerose cellule mucipare caliciformi.

% ---------------------------------------------------------------------------
\subsection{Epiteli stratificati (pluristratificati)}

\rubrica{Epitelio pavimentoso stratificato}
Distinto in \textbf{non cheratinizzato} (es.\ cervice uterina/esocervice) e \textbf{cheratinizzato} (\textbf{epidermide}).
\begin{itemize}
  \item epitelio pluristratificato piatto non cheratinizzato: riveste lo strato più esterno di bocca e vagina, oltre all'esocervice; protegge da abrasioni, disidratazioni e infezioni;
  \item epitelio pluristratificato piatto cheratinizzato: costituisce l'epidermide.
\end{itemize}

\rubrica{Epidermide: cheratinizzazione}
Le cellule (\textbf{cheratinociti}) $\rightarrow$ cheratinizzazione $\rightarrow$ il citoplasma si riempie di filamenti di cheratina $\rightarrow$ le cellule muoiono trasformandosi in lamelle cornee desquamanti. Le lamelle che costituiscono lo strato corneo più superficiale desquamano continuamente e sono rimpiazzate dalle cellule degli strati più profondi. Solo le cellule che poggiano sulla lamina basale sono in grado di proliferare, perché ricevono stimoli attivatori della proliferazione dalla matrice connettivale sottostante.

\rubrica{Epidermide: strati}
\begin{itemize}
  \item \textbf{strato basale}: cellule cubiche o cilindriche con intensa attività proliferativa;
  \item \textbf{strato spinoso}: cellule di forma poliedrica leggermente appiattita, 3--7 file di cellule, con brevi prolungamenti del citoplasma (\textit{spine});
  \item \textbf{strato granuloso}: cellule appiattite e allungate, 2--6 file di cellule, con evidenti alterazioni apoptotiche;
  \item \textbf{strato lucido}: una o più file di cellule appiattite e allungate, con filamenti di cheratina impacchettati e paralleli alla superficie della cute;
  \item \textbf{strato corneo}: lamelle cornee, presenza soltanto di cheratina, mancano nucleo e organuli cellulari.
\end{itemize}

\rubrica{Altre cellule dell'epidermide}
\begin{itemize}
  \item \textbf{melanociti}: forma stellata, producono melanina;
  \item \textbf{cellule di Langerhans}: forma stellata, presenti nello strato basale e spinoso, captano gli antigeni che attraversano l'epidermide e migrano nei linfonodi, dove inizia la risposta immune;
  \item \textbf{cellule di Merkel}: nello strato basale dell'epidermide, sono elementi sensoriali.
\end{itemize}

\rubrica{Epitelio cubico stratificato}
Riveste dotti escretori di grosso calibro (es.\ ghiandole salivari maggiori).

\rubrica{Epitelio cilindrico stratificato}
Riveste dotti escretori di grosso calibro (es.\ ghiandola sottomandibolare), spesso provvisto di ciglia sulla superficie libera.

% ---------------------------------------------------------------------------
\subsection{Epitelio di transizione}
\label{ssec:transizione}

Detto anche \textbf{urotelio}. Composto da cellule cubiche pluristratificate; le cellule più superficiali, sotto tensione, si appiattiscono (\textit{cellule ad ombrello}). Riveste la \textbf{vescica}: quando la parete dell'organo è rilasciata l'epitelio appare più alto e stratificato, quando la parete è distesa le cellule si appiattiscono.

% ---------------------------------------------------------------------------
\subsection{Specializzazioni della membrana plasmatica}

\begin{itemize}
  \item \textbf{microvilli}: estroflessioni digitiformi del citoplasma con funzione assorbente; diametro 50--100\,nm, altezza 1--2\,$\mu$m; aumentano la superficie assorbente di 25--50 volte per ogni cellula;
  \item \textbf{ciglia vibratili}: presenti nelle vie aeree, dove convogliano verso l'esterno particelle solide; presenti nelle vie genitali femminili, dove facilitano la progressione della cellula uovo dalla tuba verso l'utero;
  \item \textbf{stereociglia}: presenti nell'orecchio interno, strutture ricche in actina che vibrano in risposta alle onde sonore; presenti nel canale dell'epididimo e nel dotto deferente, dove hanno un ruolo nei processi di secrezione e riassorbimento del liquido prodotto nei tubuli seminiferi.
\end{itemize}

% ---------------------------------------------------------------------------
\subsection{Connessioni intercellulari tra cellule epiteliali}

\begin{itemize}
  \item \textbf{giunzioni occludenti}: avvolgono circolarmente l'intera cellula, impediscono il passaggio di materiale tra gli interstizi cellulari;
  \item \textbf{giunzioni aderenti}: al di sotto delle giunzioni occludenti, le E-caderine sono saldate tra loro nello spazio intercellulare e ai filamenti di actina;
  \item \textbf{desmosomi}: le E-caderine sono associate alla placca, i filamenti intermedi formano ripiegamenti a forcina;
  \item \textbf{giunzioni comunicanti}: zone di comunicazione intercellulare, permettono il trasporto di ioni e piccole molecole tra cellule;
  \item \textbf{emidesmosomi}: servono per l'adesione della parte basale della cellula epiteliale alla sottostante membrana basale.
\end{itemize}

% ---------------------------------------------------------------------------
\subsection{Epitelio ghiandolare}

\textbf{Epitelio ghiandolare}: epitelio che produce (\textit{secrezione}) sostanze che vengono espulse dalle cellule (\textit{secreto}), organizzate in strutture più o meno complesse (\textit{ghiandola}).

\begin{center}
\begin{forest} classalbero
[Ghiandole
  [Esocrine [riversano i secreti sulla superficie degli epiteli attraverso dotti escretori]]
  [Endocrine [riversano i secreti (ormoni) nel liquido interstiziale o direttamente nel sangue]]
]
\end{forest}
\end{center}

% ---------------------------------------------------------------------------
\subsection{Ghiandole esocrine}

\rubrica{Ghiandole unicellulari}
Es.\ cellula caliciforme mucipara dell'epitelio cilindrico semplice dell'intestino tenue: il nucleo è in posizione basale, mentre i numerosi granuli di secrezione occupano la regione apicale espansa della cellula.

\rubrica{Tipologia del secreto}
\begin{itemize}
  \item \textbf{ghiandola sierosa}: il secreto è una soluzione acquosa di una o più sostanze (es.\ saliva);
  \item \textbf{ghiandola mucosa}: il secreto è formato da mucine (es.\ succo duodenale);
  \item \textbf{ghiandola mista}: adenomero misto, con secreto sia sieroso sia mucoso (es.\ semiluna sierosa associata ad acino mucoso).
\end{itemize}

\rubrica{Criteri di classificazione delle ghiandole esocrine pluricellulari}
\begin{itemize}
  \item in base al \textbf{dotto escretore}: ghiandola semplice (dotto non ramificato), ramificata (dotto non ramificato ma adenomero ramificato), composta (dotto ramificato);
  \item in base alla \textbf{forma dell'adenomero}: ghiandola tubulare, acinosa, alveolare;
  \item in base alla \textbf{modalità di secrezione}: ghiandola olocrina, apocrina, merocrina.
\end{itemize}

\rubrica{Modalità di secrezione}
\begin{itemize}
  \item \textbf{merocrina} (o \textit{eccrina}): secrezione attraverso vescicole, regolata da un normale meccanismo di esocitosi, senza danni al citoplasma (es.\ muco, succhi gastrici, pancreatici, intestinali, bile); è il tipo di secrezione più comune;
  \item \textbf{apocrina}: i granuli di secreto si accumulano nella parte apicale della cellula, che al momento della secrezione si disgrega ed entra a far parte del secreto, con perdita di parte del citoplasma (es.\ latte);
  \item \textbf{olocrina}: all'atto della secrezione si ha il disfacimento dell'intera cellula, con distruzione completa della cellula (es.\ sebo).
\end{itemize}

% ---------------------------------------------------------------------------
\subsection{Ghiandole endocrine}

Riversano i loro secreti (ormoni) nel liquido intercellulare interstiziale o direttamente nel sangue: il messaggero chimico, destinato ad agire su cellule bersaglio distanti dal luogo della sua sintesi, viene immesso nel torrente circolatorio; per questo motivo le cellule a secrezione endocrina sono sempre localizzate nelle vicinanze di capillari sanguigni.

\rubrica{Organizzazione strutturale}
\begin{itemize}
  \item \textbf{follicoli}: es.\ follicolo tiroideo, rivestito da un epitelio cubico o cilindrico monostratificato, contenente la \textit{colloide} (deposito degli ormoni tiroidei); le cellule follicolari costituiscono l'epitelio, mentre le \textbf{cellule parafollicolari} sono poste nel connettivo circostante e producono calcitonina;
  \item \textbf{cordoni};
  \item \textbf{nidi}.
\end{itemize}

\rubrica{Meccanismo d'azione degli ormoni}
\begin{itemize}
  \item \textbf{ormoni peptidici}: (1) l'ormone lega un recettore proteico di membrana; (2) la proteina recettoriale attiva una via di segnalazione nella cellula; (3) una serie di molecole intermedie (\textit{trasduttori}) trasmette il segnale ad una proteina che esegue la risposta cellulare (es.\ scissione del glicogeno in glucosio);
  \item \textbf{ormoni steroidei}: (1) l'ormone entra nella cellula per diffusione attraverso la membrana; (2) si lega a un recettore intracellulare; (3) il complesso ormone-recettore penetra nel nucleo; (4) il complesso ormone-recettore si lega al DNA, accendendo o spegnendo la trascrizione di specifici geni $\rightarrow$ risposta cellulare: attivazione di un gene e sintesi di una nuova proteina.
\end{itemize}
